O algoritmo Insertion Sort apresenta diferentes comportamentos de
desempenho dependendo da organização inicial dos elementos do vetor.
Sua complexidade de tempo é analisada considerando o melhor caso, o caso
médio e o pior caso.

\subsubsection{Melhor caso -- \texorpdfstring{$\Theta(n)$}{Theta(n)}}

O melhor caso ocorre quando o vetor já está ordenado de forma crescente.
Nessa situação, para cada elemento analisado, o algoritmo realiza apenas
uma comparação com o elemento anterior e verifica que não é necessário
deslocar nenhum elemento.

Considerando um vetor com $n$ elementos, o laço externo é executado
$n-1$ vezes. Em cada iteração, é realizada uma quantidade constante de
operações, como atribuições e comparações.

Assim, o número total de operações é proporcional a:

\begin{equation}
T(n) = c(n-1) + d
\end{equation}

em que $c$ e $d$ são constantes. Desconsiderando os termos constantes e
os de menor crescimento, temos:

\begin{equation}
T(n) = \Theta(n)
\end{equation}

Portanto, o melhor caso do Insertion Sort possui complexidade linear.

\subsubsection{Pior caso -- \texorpdfstring{$\Theta(n^2)$}{Theta(n²)}}

O pior caso ocorre quando o vetor está ordenado de forma decrescente e
deseja-se ordená-lo de forma crescente. Nesse cenário, cada novo elemento
precisa ser comparado com todos os elementos anteriores e deslocado até o
início da parte já ordenada.

A quantidade total de deslocamentos pode ser representada pela soma:

\begin{equation}
T(n) = 1 + 2 + 3 + \cdots + (n-1)
\end{equation}

Essa soma corresponde à soma dos termos de uma progressão aritmética:

\begin{equation}
T(n) = \frac{n(n-1)}{2}
\end{equation}

Desenvolvendo a expressão:

\begin{equation}
T(n) = \frac{n^2-n}{2}
\end{equation}

O termo de maior crescimento é $n^2$. Portanto:

\begin{equation}
T(n) = \Theta(n^2)
\end{equation}

Assim, o pior caso do Insertion Sort apresenta complexidade quadrática.

\subsubsection{Caso médio -- \texorpdfstring{$\Theta(n^2)$}{Theta(n²)}}

O caso médio ocorre quando os elementos estão distribuídos aleatoriamente.
Nessa situação, o algoritmo realiza, em média, aproximadamente metade dos
deslocamentos possíveis para cada elemento inserido.

Temos:

\begin{equation}
T(n) \approx \frac{1}{2}
\left(1+2+3+\cdots+(n-1)\right)
\end{equation}

Aplicando a fórmula da soma:

\begin{equation}
T(n) \approx \frac{1}{2}
\cdot\frac{n(n-1)}{2}
\end{equation}

\begin{equation}
T(n) \approx \frac{n^2-n}{4}
\end{equation}

Como o termo dominante é $n^2$, a complexidade assintótica é:

\begin{equation}
T(n) = \Theta(n^2)
\end{equation}

Portanto, no caso médio, o Insertion Sort também apresenta complexidade
quadrática.

\subsubsection{Resumo das complexidades}

\begin{table}[H]
    \centering
    \caption{Complexidades do algoritmo Insertion Sort}
    \label{tab:complexidade-insertion}
    \begin{tabular}{lc}
        \hline
        \textbf{Caso} & \textbf{Complexidade de tempo} \\
        \hline
        Melhor caso & $\Theta(n)$ \\
        Caso médio & $\Theta(n^2)$ \\
        Pior caso & $\Theta(n^2)$ \\
        \hline
    \end{tabular}
\end{table}

Conclui-se que o Insertion Sort é eficiente para vetores pequenos ou
parcialmente ordenados, mas apresenta desempenho inferior para grandes
conjuntos de dados desordenados.